\section{Семинар \thesection}
\subsection{Задачи по тождественным частицам.}
Выписываю две задачи из задавальника и одну мою, которые нужно порешать.
Первая это упражнение 9, вторая моя, третья задача номер 13.
Если нужно решение для моей задачи, напиши, я тебе его скину сразу
(вдруг ты хочешь сама посчитать хз).
\begin{exercise}
  \textit{Используя понижающий оператор $\hat{S}_-$, постройте
  собственные функции операторов $\hat{\mathbf{S}}^2$ и $\hat{S}_z$
  двух бозонов со спинами s_1 = s_2 = 1
  ($\hat{\mathbf{S}} = \hat{s}_1 + \hat{s}_2$, $\hat{S}_- = \hat{s}_{1-} + \hat{s}_{2-}$)}
\end{exercise}
\qedsq


\begin{exercise}
  \textit{Посчитайте коэффициенты Клебша-Гордона для задачи сложения двух моментов:
  спиновый момент электрона (спин 1/2) и момент орбиты ($j_1 = l$, где $l \in \mathbb{Z}$).}
\end{exercise}
\qedsq

\begin{exercise}
  \textit{Найдите уровни энергии и волновые функции стационарных состояний двух
  невзаимодействующих тождественных частиц в потенциальном ящике
  \[
  V(x) = \begin{cases} 0, & 0 < x < a,\\ \infty, & x<0,\, x>a, \end{cases}
  \]
  если этими частицами являются: (а) ферми-частицы со спинами $s=1/2$;
  (б) бозе-частицы со спинами s=0; (в) бозе-частицы со спинами s=1.
  Чему равна в каждом из этих случаев энергия основного состояния N частиц?
  }
\end{exercise}
\qedsq

Если так получится, что семинар выйдет не прям большим (4-6 страниц), можно
добавить ещё одну задачу на тождественные частицы. Я в целом затрагиваю
гелий в семинаре, поэтому можно взять задачу 5.15 или 5.16 из гриффита
(какая тебе кажется более интересной). Только не бери пж ту, которая очевидно проще.
Если будет откровенно лень писать, я лучше сам напишу
